\documentclass[12pt,a4paper,oneside]{article}
\usepackage[brazil]{babel}
\usepackage[utf8]{inputenc}
\usepackage{olymp}
\usepackage{tikz}

\setcounter{problem}{3}

\begin{document}

\begin{problem}{Montes}{montes}{2}
 
Na próxima semana será realizada em Montes Claros a VII Maratona Mineira de Programação e três equipes de alunos da UFV vão participar, inclusive alunos desta turma de INF110. Para passar o tempo nas mais de 9 horas de viagem, resolveram anotar a altura do relevo da estrada em intervalos regulares. A ideia é que, quando voltarem, possam determinar o número de montes e de picos encontrados ao longo da viagem. 

Considere, como exemplo, que anotaram os seguintes valores:

7 5 8 7 3 0 0 0 4 8 10 8 6 3 0 0 4 10 14 11 12 8 7 0 0

Um desenho mostrando o relevo representado por essas anotações é mostrado na figura abaixo. Note que existem 3 montes (coloridos de cinza) e 4 picos (pontos destacados).

\includegraphics[scale=0.5]{images/exercise_106_0}
% \begin{tikzpicture}[scale=0.4, transform shape]
% \draw[gray, thick] (0,7) -- (1,5) -- (2,8) -- (3,7) -- (4,3) -- (5,0) -- (6,0) -- (7,0) -- (8,4) -- (9,8) -- (10,10) -- (11,8) -- (12,6) -- (13,3) -- (14,0) -- (15,0) -- (16,4) -- (17,10) -- (18,14) -- (19,11) -- (20,12) -- (21,8) -- (22,7) -- (23,0) -- (24,0)
% \end{tikzpicture}

Os montes são separados por pontos ou regiões de planície (altura 0). E os picos são pontos de máximo local, ou seja, de altura maior que as alturas anotadas imediatamente antes e depois.

 A primeira e a última altura anotadas nunca são considerados picos, pois não se sabe o que vem antes ou depois (veja que no exemplo dado acima o primeiro ponto não foi considerado pico).

\begin{minipage}{.7\textwidth}
\Notes
Não existem alturas iguais consecutivas, exceto nas planícies, que podem conter um ou mais zeros consecutivos. Assim, não existirá, por exemplo, o caso 0 3 5 5 2 0, ilustrado ao lado.
\end{minipage}\hspace{30pt}
\begin{minipage}{.2\textwidth}
\includegraphics[scale=0.4]{images/exercise_106_1}
% \begin{tikzpicture}[scale=0.4, transform shape]
% \draw[gray, thick] (0,0) -- (1,3) -- (2,5) -- (3,5) -- (4,2) -- (5,0)
% \end{tikzpicture}
\end{minipage}

\InputFile

A entrada é composta por duas linhas. A primeira contém um valor inteiro $N$, que indica o número de anotações tomadas ao longo da viagem. A segunda linha contém $N$ valores inteiros, $A_i$, indicando a altura anotada na $N$-ésima anotação. Restrições: $1 \leq N \leq 10000$; $0 \leq A_i \leq 1000$.

\OutputFile

Seu programa deve imprimir dois valores inteiros separados por espaço em branco: o número de montes e o número de picos.

% \Notes
% Preste atenção aos tipos das variáveis, pois $F(90) \approx 2^{61}$.

\Example

\begin{examplewide}
\exmp{
24
7 5 8 7 3 0 0 0 4 8 10 8 6 3 0 0 4 10 14 11 12 8 7 0 0
}{
3 4
}%
\end{examplewide}

\begin{examplewide}
\exmp{
8
1 0 1 0 1 0 1 0
}{
4 3
}%
\end{examplewide}

\begin{examplewide}
\exmp{
1
5
}{
1 0
}%
\end{examplewide}

\begin{examplewide}
\exmp{
2
0 8
}{
1 0
}%
\end{examplewide}

\begin{examplewide}
\exmp{
1
0
}{
0 0
}%
\end{examplewide}

\begin{examplewide}
\exmp{
13
0 1 2 3 2 1 0 1 2 3 2 1 0
}{
2 2
}%
\end{examplewide}


\begin{examplewide}
\exmp{
4
10 7 5 2
}{
1 0
}%
\end{examplewide}

%Comentários sobre o exemplo, se necessário.

\end{problem}

\end{document}
